\clearpage
\fancyhead[L]{\small VESTA.}
\hypertarget{native-293-lower}{}
\begin{center}
VESTA.
\end{center}

\begin{center}
\rule{0.6\linewidth}{0.4pt}\\
Monatliche Correspondenz. Band XVI. S. 83–92. 1807 Juli.\\
\rule{0.6\linewidth}{0.4pt}
\end{center}

\par\indent Kaum ist seit der ersten zum drittenmal einem Deutschen gelungenen Entdeckung eines neuen Planeten, der \emph{Vesta}, ein Zeitraum von drei Monaten verflossen, und schon wurde dieses kleine Gestirn auf den meisten in- und ausländischen Sternwarten aufgefunden, beobachtet, und so dem Geometer die Data geliefert, um die Bahn des neuen Irrsterns zu erforschen und in engere Grenzen einzuschliessen. Wenn die schnelle Vermehrung dieser schwierigen Beobachtungen, wenn die systematische nicht zufällige Art der Entdeckung dieses neuen Planeten eine neue merkwürdige glänzende Epoche in dem praktischen Zustande der beobachtenden Astronomie bezeichnet, so muss es auf der andern Seite den Mathematiker, der bekannt ist mit frühern Bemühungen über Methoden zu Bestimmung von Planeten-Bahnen, bekannt mit den weitläuftig ermüdenden Rechnungen, die alle zeither gegebene directe und indirecte Verfahrungs-Arten erforderten, in ein freudiges Erstaunen versetzen, wenn er sieht, wie die ersten Elemente dieses Planeten, nach einer nur 10stündigen Arbeit, schon als sehr genähert, und bald darauf die zweiten, aus einem durchlaufenen kleinen heliocentrischen Bogen von noch nicht $11^\circ$ mit einer Sicherheit und Schärfe dargestellt werden, dass die nach diesen II. Elementen berechneten Örter der \emph{Vesta} mit 84 beobachteten, auf eine Art harmoniren, wie es kaum bei ältern Planeten-Tafeln der Fall immer ist.

\par\indent Meine zweiten Elemente der \emph{Vesta}, schreibt uns unterm 7. Juni der Geometer, dem Bestimmung von Planeten-Bahnen nur leichtes Spiel zu sein scheint, die ich Ihnen vor einiger Zeit überschickte, haben sich so brav gehalten, dass ich bisher noch keine Veranlassung gehabt habe, etwas daran zu verbessern. Die Beobachtungen seit der Mitte des Mai zeigen zwar, dass die Elemente die Rectascensionen etwas zu gross geben, allein diese Differenz ist noch in sehr engen Grenzen geblieben. Die Declinationen stimmen fortdauernd vortrefflich. Wahrscheinlich werden die Elemente bis zu dem Zeitpunkte, wo die \emph{Vesta} für dieses Jahr unsichtbar werden wird, sich um etwa $1'$ von den Beobachtungen entfernen können, daher ich es für überflüssig halte, früher eine neue Bestimmung der Bahn vorzunehmen. Für diesmal schicke ich Ihnen also nur die sorgfältig gemachte Vergleichung aller mir bekannt gewordenen Beobachtungen mit den zweiten Elementen; die Differenzen bis etwa zum 6. Mai wird man dreist dem grössern Theile nach auf die respectiven Beobachtungen schieben können. Die meisten dieser Beobachtungen (34) sind zu \emph{Bremen} angestellt, 15 zu \emph{Göttingen}, 12 zu \emph{Lilienthal}, 10 in \emph{Berlin}, 5 in \emph{Greenwich}, 3 in \emph{Paris}, 3 in \emph{Prag}, 2 in \emph{Braunschweig}, zusammen 84.
\clearpage
\begingroup
\fontsize{9}{11}\selectfont
\setlength{\tabcolsep}{2.5pt}\renewcommand{\arraystretch}{1.10}
\begin{longtable}{lrrrrrl}
1807 & \shortstack{Mittlere Zeit\\am Beob.-Orte} & \multicolumn{2}{c}{Berechnete} & \multicolumn{2}{c}{Unterschied R.$-$B.} & Beobachter. \\
& & AR. & Decl. & AR. & Decl. & \\ \hline
\SourceLeaf{294}März 29 & $8^{\text{u}}\;21^{\text{m}}\;20^{\text{s}}$ & $184^\circ\,8'\,53''4$ &  & $-2''4$ &  & Olbers \\
 & $10\;31\;16$ & $184\;7\;39.1$ & $11^\circ 47' 56''0$ & $-11.7$ & $+9''0$ &  \\
30 & $8\;34\;53$ & $183\;55\;\;1.8$ &  & $+5.0$ &  &  \\
 & $8\;44\;8$ & $183\;54\;56.5$ & $11\;53\;21.7$ & $+10.7$ & $+10.7$ &  \\
 & $12\;33\;17$ & $183\;52\;45.6$ & $11\;54\;16.5$ & $+4.8$ & $-10.5$ &  \\
April 1 & $9\;50\;0$ & $183\;27\;\;2.8$ & $12\;\;4\;41.6$ & $0$ & $+10.4$ &  \\
 & $12\;4\;35$ & $183\;25\;47.1$ & $12\;\;5\;11.2$ & $+22.9$ & $+20.5$ & Bessel \\
2 & $8\;21\;1$ & $183\;14\;23.2$ & $12\;\;9\;34.7$ & $+5.4$ & $-12.3$ & Olbers \\
 & $11\;38\;13$ & $183\;12\;33.1$ & $12\;10\;16.3$ & $+19.0$ & $-17.9$ & Bessel \\
3 & $8\;16\;49$ & $183\;\;1\;\;3.3$ & $12\;14\;32.5$ & $+2.5$ & $-7.5$ & Olbers \\
4 & $9\;9\;43$ & $182\;47\;20.8$ & $12\;19\;27.2$ & $0$ & $-17.8$ &  \\
5 & $11\;17\;3$ & $182\;33\;11.0$ & $12\;24\;18.9$ & $+0.3$ & $-0.2$ & Harding \\
6 & $8\;22\;45$ & $182\;21\;54.0$ & $12\;28\;\;1.4$ & $-0.6$ & $+0.5$ & Gauss \\
 & $11\;12\;15$ & $182\;20\;22.4$ & $12\;28\;31.1$ & $-11.5$ & $+41.7$ & Harding \\
8 & $8\;21\;37$ & $181\;56\;44.8$ & $12\;35\;45.4$ & $-6.1$ & $-0.4$ & Gauss \\
 & $8\;28\;49$ & $181\;56\;44.5$ & $12\;35\;45.5$ & $+1.5$ & $-12.5$ & Olbers \\
 & $11\;2\;43$ & $181\;55\;24.2$ & $12\;36\;\;9.0$ & $+5.0$ & $-3.2$ & Harding \\
9 & $8\;30\;46$ & $181\;44\;27.2$ & $12\;39\;15.0$ & $+10.2$ & $+1.0$ & Olbers \\
 & $10\;57\;58$ & $181\;43\;16.0$ & $12\;39\;34.7$ & $+5.0$ & $-1.4$ & Harding \\
12 & $8\;27\;20$ & $181\;\;9\;28.1$ & $12\;48\;\;2.0$ & $-5\;9$ & $-14.0$ & Olbers \\
13 & $10\;39\;14$ & $180\;57\;31\;2$ & $12\;50\;36.1$ & $+14.2$ & $-6.9$ & Bode \\
14 & $8\;19\;51$ & $180\;47\;39.9$ & $12\;52\;32.2$ & $-18.1$ & $-15.8$ & Olbers \\
 & $10\;43\;33$ & $180\;46\;42.5$ & $12\;52\;42.7$ & $+2.6$ & $+3.0$ & Harding \\
 & $10\;34\;36$ & $180\;46\;29.1$ & $12\;52\;45.2$ & $+3.9$ & $+7.0$ & Burckhardt \\
17 & $8\;12\;35$ & $180\;17\;24.9$ & $12\;57\;14.7$ & $-1.9$ & $-1.3$ & Olbers \\
23 & $8\;30\;52$ & $179\;26\;39.4$ & $12\;59\;23.6$ & $-1.6$ & $+3.6$ &  \\
 & $9\;53\;49$ & $179\;26\;15.7$ & $12\;59\;22.1$ & $-1.4$ & $+7.8$ & Harding \\
24 & $9\;49\;26$ & $179\;19\;19.2$ & $12\;58\;47.2$ & $+3\;2$ & $-7.8$ & Bode \\
25 & $8\;44\;38$ & $179\;12\;57.9$ & $12\;57\;58.8$ & $-21.1$ & $+0.8$ & Olbers \\
 &  & $179\;12\;54.4$ & $12\;57\;58.2$ & $+5.4$ & $-7.8$ &  \\
 & $9\;45\;5$ & $179\;12\;43.1$ & $12\;57\;56.5$ & $+7.1$ & $-10.5$ & Bode \\
 &  & $179\;12\;28.9$ & $12\;57\;54.3$ & $-1.1$ & $+8.3$ & Groombridge \\
 & $10\;25\;19$ & $179\;12\;27.1$ & $12\;57\;54.1$ & $+4.0$ & $+8.3$ & Bessel \\
26 & $8\;46\;4$ & $179\;\;6\;41.1$ & $12\;56\;52.4$ & $-6.7$ & $+1.4$ & Olbers \\
 & $9\;40\;45$ & $179\;\;6\;31.7$ & $12\;56\;50.5$ & $-4.3$ & $-7.5$ & Bode \\
 & $9\;40\;42$ & $179\;\;6\;28.3$ & $12\;56\;49.8$ & $+8.5$ & $+1.1$ & Harding \\
 & $9\;40\;45$ & $179\;\;6\;20.7$ & $12\;56\;48.3$ & $+3.2$ & $2.9$ & Burckhardt \\
 &  & $179\;\;6\;18.4$ & $12\;56\;47.8$ & $-6.6$ & $+3.8$ & Groombridge \\
 & $10\;44\;12$ & $179\;\;6\;11.4$ & $12\;56\;46.3$ & $+0.7$ & $+2.0$ & Bessel \\
27 & $9\;36\;26$ & $179\;\;0\;45.6$ & $12\;55\;29\;4$ & $-2.4$ & $+7.4$ & Bode \\
 & $9\;36\;23$ & $179\;\;0\;42.3$ & $12\;55\;28.6$ & $+10.2$ & $-0\;1$ & Harding \\
 &  & $179\;\;0\;33.1$ & $12\;55\;26.1$ & $+1.1$ & $-3.9$ & Maskelyne \\
 & $10\;45\;27$ & $179\;\;0\;25.4$ & $12\;55\;24.0$ & $-1.4$ & $-0.2$ & Bessel \\
 & $11\;18\;35$ & $179\;\;0\;17.7$ & $12\;55\;21.9$ & $+6.7$ & $+0.9$ & Olbers \\
28 &  & $178\;55\;15.2$ & $12\;53\;49.5$ & $+6.2$ & $+5.5$ & Groombridge \\
 & $10\;57\;6$ & $178\;55\;\;2.6$ & $12\;53\;45.9$ & $+6.4$ & $0$ & Bessel \\
 & $11\;34\;20$ & $178\;54\;54.4$ & $12\;53\;43.1$ & $-4.6$ & $-7.8$ & Olbers \\
29 & $9\;27\;52$ & $178\;50\;29.2$ & $12\;52\;\;2.4$ & $+8.2$ & $+0.4$ & Bode \\
 & $9\;28\;51$ & $178\;50\;26.5$ & $12\;52\;\;1.2$ & $-13.5$ & $+2.8$ & Harding \\
 &  & $178\;50\;18.8$ & $12\;51\;58.0$ & $+3.8$ & $+6.0$ & Maskelyne \\
 & $11\;4\;54$ & $178\;50\;\;6.7$ & $12\;51\;52.9$ & $-11.9$ & $-7.7$ & Bessel \\
30 & $9\;23\;38$ & $178\;45\;59.6$ & $12\;49\;56\;7$ & $+11.6$ & $-3.3$ & Bode \\
 & $9\;23\;39$ & $178\;45\;57.2$ & $12\;49\;55.4$ & $-2.9$ & $-5.2$ & Harding \\
 & $11\;27\;45$ & $178\;45\;34.4$ & $12\;49\;43.5$ & $+2.2$ & $-9.6$ & Bessel \\
Mai 1 & $9\;19\;26$ & $178\;41\;56.4$ & $12\;47\;36.5$ & $+0.4$ & $+5.5$ & Bode \\
 & $11\;23\;35$ & $178\;41\;33.6$ & $12\;47\;21.8$ & $+7.6$ & $-30.2$ & Olbers \\
 & $11\;57\;53$ & $178\;41\;28.3$ & $12\;47\;18.3$ & $+12.0$ & $-6.7$ & Bessel \\
2 & $9\;15\;14$ & $178\;38\;17.3$ & $12\;45\;\;0.3$ & $+7.1$ & $-3.8$ & Harding \\
\end{longtable}
\endgroup

\clearpage
\begingroup
\fontsize{9}{11}\selectfont
\setlength{\tabcolsep}{2.5pt}\renewcommand{\arraystretch}{1.10}
\begin{longtable}{lrrrrrl}
1807 & \shortstack{Mittlere Zeit\\am Beob.-Orte} & \multicolumn{2}{c}{Berechnete} & \multicolumn{2}{c}{Unterschied R$-$B.} & Beobachter. \\
& & AR. & Decl. & AR. & Decl. & \\ \hline
\SourceLeaf{295}Mai 4 & $9^{\text{u}}\;5^{\text{m}}\;43^{\text{s}}$ & $178^\circ 32'\,24''4$ & $12^\circ 39' 10''8$ & $-0''6$ & $-20''2$ & David \\
 & $9\;7\;1$ & $178\;32\;23.9$ & $12\;39\;10.2$ & $-5.1$ & $+6.2$ & Bode \\
 & $9\;6\;59$ & $178\;32\;22.2$ & $12\;39\;\;8.0$ & $-1.7$ & $+3.2$ & Harding \\
5 & $9\;1\;37$ & $178\;30\;\;6.6$ & $12\;35\;54.2$ & $+4.6$ & $-11.8$ & David \\
 & $9\;2\;56$ & $178\;30\;\;6.1$ & $12\;35\;53.4$ & $+10.1$ & $+4.4$ & Bode \\
 & $9\;2\;54$ & $178\;30\;\;4.9$ & $12\;35\;51\;5$ & $+1.7$ & $-4.2$ & Harding \\
 & $9\;4\;36$ & $178\;30\;\;4.4$ & $12\;35\;50.6$ & $+7.4$ & $-2.4$ & Olbers \\
 & $11\;33\;55$ & $178\;29\;51.6$ & $12\;35\;29.3$ & $+1.6$ & $-1.7$ &  \\
6 & $8\;57\;32$ & $178\;28\;15.3$ & $12\;32\;23.8$ & $+20.3$ & $-10.2$ & David \\
 & $9\;16\;10$ & $178\;28\;12.6$ & $12\;32\;17.7$ & $+3.6$ & $-5.3$ & Olbers \\
8 & $9\;51\;51$ & $178\;25\;51.7$ & $12\;24\;39.9$ & $+4.6$ & $+1.6$ & Harding \\
 & $9\;18\;20$ & $178\;25\;50.7$ & $12\;24\;34.6$ & $+1.7$ & $+3.6$ & Olbers \\
 & $10\;30\;36$ & $178\;25\;48.5$ & $12\;24\;22.4$ & $+11.6$ & $-12.7$ & Bessel \\
11 &  & $178\;25\;37.4$ & $12\;11\;32.2$ & $+2.2$ & $+5.2$ & Harding \\
 & $10\;31\;39$ & $178\;25\;40.3$ & $12\;11\;\;9.2$ & $+3.5$ & $-2.0$ & Bessel \\
15 & $8\;23\;43$ & $178\;31\;26.4$ & $11\;51\;\;4.9$ & $+9\;0$ & $-2.8$ & Burckhardt \\
19 & $10\;43\;2$ & $178\;44\;20.4$ & $11\;27\;10.4$ & $-4.4$ & $-13.1$ & Bessel \\
 & $11\;28\;15$ & $178\;44\;28.0$ & $11\;26\;59.7$ & $+25.0$ & $-71.3$ & Olbers \\
22 & $10\;23\;12$ & $178\;58\;\;4.5$ &  & $+17.5$ &  &  \\
23 & $10\;52\;10$ & $179\;\;3\;43.2$ & $11\;\;0\;55.1$ & $+17.2$ & $-10.9$ &  \\
24 & $11\;28\;23$ & $179\;\;9\;29.8$ & $10\;53\;47.9$ & $+17.8$ & $+2.9$ &  \\
25 & $10\;29\;53$ & $179\;15\;24.3$ & $10\;46\;58.5$ & $+19.6$ & $-1.6$ &  \\
 & $10\;53\;31$ & $179\;15\;30.5$ & $10\;46\;51.5$ & $+17.7$ & $+2.9$ &  \\
26 & $10\;45\;20$ & $179\;22\;\;0.2$ & $10\;39\;38.4$ & $+29.0$ & $-6.9$ &  \\
31 & $11\;32\;54$ & $180\;\;0\;25.9$ &  & $+22.9$ &  &  \\
Juni 2 & $12\;19\;7$ & $180\;18\;28.9$ & $9\;44\;18.4$ & $+23.9$ & $+0.4$ &  \\
\end{longtable}
\endgroup
\smallskip
\par\indent „\textsc{Harding}'s Declination vom 6. April und \textsc{Olbers} Declinationen vom 1. und 19. Mai
scheinen fehlerhaft zu sein. Bei verschiedenen Beobachtungen von \textsc{Olbers}, \textsc{Harding} und
\textsc{Bessel} ist die Vergleichung nach einigen spätern Berichtigungen gemacht, daher die Ihnen vielleicht
früher mitgetheilten Angaben etwas differiren werden.“ — [\textsc{Groombridge} und \textsc{Maskeline}
machten Meridian-Beobachtungen.]

\smallskip
\noindent---\ ---\ ---

\begin{flushright}
[\textsc{von Zach}.]
\end{flushright}
